% !TEX root = ../thesis.tex
%=========================================================
\chapter{Supplementary: Tunnelbarrier}
\label{appendix:tb}
%=========================================================

    %=========================================================
    \section{Rectification and Asymmetry}
    \label{appendix:tb:asym}
    %=========================================================

    For completeness, Fig.~\ref{fig:tb:asym:comparison-cuts} shows branch-resolved line cuts corresponding to the comparison maps in Fig.~\ref{fig:tb:asym:comparison-maps}. Figures~\ref{fig:tb:asym:sequence-iv}, \ref{fig:tb:asym:sequence-didv}, and \ref{fig:tb:asym:sequence-dvdi} show the selected frequency sweeps in the three transport representations. Figures~\ref{fig:tb:asym:nu_7.97GHz:cal} and \ref{fig:tb:asym:nu_7.97GHz:sim} provide the corresponding detailed measured and simulated transport landscapes.

    \begin{figure}[t]
        \centering
        \import{tunnelbarrier/asym/irradiation/comp-cuts}{main.pgf}
        \caption{
            Comparison of the positive- and negative-bias transport branches of TB2 at selected frequencies and calibrated microwave amplitudes. The columns correspond to the six frequencies shown in Fig.~\ref{fig:tb:asym:comparison-maps}, and the rows show the normalized \textit{I--V}, \textit{dI--dV}, and \textit{dV--dI} responses. Solid lines denote the positive-bias branch, while dashed lines denote the negative-bias branch reflected about the origin. The color encodes the normalized calibrated amplitude $eA_\mathrm{cal}/h\nu$. The increasing separation of the two branches at intermediate frequencies reveals the onset and frequency-dependent reversal of the rectification.
        }
        \label{fig:tb:asym:comparison-cuts}
    \end{figure}


    \begin{figure}[p]
        \centering
        \landscapepanels{tunnelbarrier/asym/irradiation/nu_7.97GHz}{cal}
        \caption{
            Calibrated transport landscape of TB2 at $\nu=7.97\,\mathrm{GHz}$. The three columns show the \textit{I--V}, \textit{dI--dV}, and \textit{dV--dI} responses as selected traces, surfaces, and false-color maps. The vertical axis uses the effective calibrated amplitude defined in Eq.~\eqref{eq:tb:asym:calibrated-amplitude}.
            }
            \label{fig:tb:asym:nu_7.97GHz:cal}

            \vspace{1em}

            \landscapepanels{tunnelbarrier/asym/irradiation/nu_7.97GHz}{sim}
            \caption{
            Symmetric SIS plus Tien--Gordon reference calculation at $\nu=7.97\,\mathrm{GHz}$. The simulation shows the response expected from bias-symmetric photon-assisted quasiparticle tunneling and provides a model reference for the asymmetric measurement. The temperature over the amplitude axis was calculated with the fit aquired in \ref{fig:tb:asym:control-Tfit_Afit}.
            }
            \label{fig:tb:asym:nu_7.97GHz:sim}
        \end{figure}



        % \landscapefigure{tunnelbarrier/asym/irradiation/nu_7.75GHz}{cal}{$7.75GHz$}{testexp1}
        % \landscapefigure{tunnelbarrier/asym/irradiation/nu_7.87GHz}{cal}{$7.87GHz$}{testexp2}
        % \landscapefigure{tunnelbarrier/asym/irradiation/nu_7.89GHz}{cal}{$7.89GHz$}{testexp3}
        % \landscapefigure{tunnelbarrier/asym/irradiation/nu_7.91GHz}{cal}{$7.91GHz$}{testexp4}
        % \landscapefigure{tunnelbarrier/asym/irradiation/nu_7.96GHz}{cal}{$7.96GHz$}{testexp5}
        % \landscapefigure{tunnelbarrier/asym/irradiation/nu_7.97GHz}{cal}{$7.97GHz$}{testexp6}
        % \landscapefigure{tunnelbarrier/asym/irradiation/nu_8.05GHz}{cal}{$8.05GHz$}{testexp7}
        % \landscapefigure{tunnelbarrier/asym/irradiation/nu_7.96GHz}{sim}{$7.96GHz$}{testsim5}
        % \landscapefigure{tunnelbarrier/asym/irradiation/nu_7.97GHz}{sim}{$7.97GHz$}{testsim6}


    %=========================================================
    \section{Coupling and Heating for TB2}
    \label{appendix:coupling-and-heating}
    %=========================================================

        I used the following diagnostic relations to compare microwave coupling and heating across the investigated frequencies:
        \begin{align}
            A_\mathrm{fit} &= \eta_{A_\mathrm{out}}^{(A_\mathrm{fit})}\, A_\mathrm{out}\,,\\
            T_\mathrm{fit} &= \min\!\left[T_\mathrm{c},\,\max\!\left(T_\mathrm{base},\,\eta_{A_\mathrm{out}}^{(T_\mathrm{fit})}\, A_\mathrm{out}+T_\mathrm{off}\right)\right]\,,\\
            T_\mathrm{bath} &= \eta_{A_\mathrm{out}}^{(T_\mathrm{bath})}\, A_\mathrm{out}+ \mathrm{const.}\,,\\
            T_\mathrm{fit} &= \min\!\left[T_\mathrm{c},\,\max\!\left(T_\mathrm{base},\,\eta_{A_\mathrm{fit}}^{(T_\mathrm{fit})}\, A_\mathrm{fit}+T_\mathrm{off}\right)\right]\,,\\
            T_\mathrm{bath} &= \eta_{A_\mathrm{fit}}^{(T_\mathrm{bath})}\, A_\mathrm{fit}+ \mathrm{const.}\,,\\
            T_\mathrm{fit} &= \min\!\left[T_\mathrm{c},\,\max\!\left(T_\mathrm{base},\,\eta_{T_\mathrm{bath}}^{(T_\mathrm{fit})}\, T_\mathrm{bath}+T_\mathrm{off}\right)\right]\,.
            \label{eq:tb:asym:diagnostic-couplings}
        \end{align}

        The fitted-temperature data formed empirical lower and upper clusters near $T_\mathrm{base}\approx0.47(1)\,\mathrm{K}$ and $T_\mathrm{c}\approx1.29(1)\,\mathrm{K}$. I used these values as the bounds in the clipped temperature relations. The coefficients $\eta$ quantify correlations between the respective coordinates and do not by themselves identify a causal coupling mechanism.

        \begin{figure}[t]
            \centering
            \import{tunnelbarrier/asym/irradiation/control}{Tfit_Aout.pgf}
            \caption{
                Fitted transport temperature $T_\mathrm{fit}$ as a function of the nominal source amplitude $A_\mathrm{out}$ at selected frequencies. The clipped fits use the empirical bounds identified from the temperature clusters.
            }
            \label{fig:tb:asym:control-Tfit_Aout}
        \end{figure}

        \begin{figure}[t]
            \centering
            \import{tunnelbarrier/asym/irradiation/control}{Tbath_Aout.pgf}
            \caption{
                Measured bath temperature $T_\mathrm{bath}$ as a function of the nominal source amplitude $A_\mathrm{out}$ at selected frequencies. The slopes quantify microwave-induced bath heating for the frequency comparison.
            }
            \label{fig:tb:asym:control-Tbath_Aout}
        \end{figure}

        \begin{figure}[t]
            \centering
            \import{tunnelbarrier/asym/irradiation/control}{Tfit_Afit.pgf}
            \caption{
                Fitted transport temperature $T_\mathrm{fit}$ as a function of the effective fitted microwave amplitude $A_\mathrm{fit}$ at selected frequencies.
            }
            \label{fig:tb:asym:control-Tfit_Afit}
        \end{figure}

        \begin{figure}[t]
            \centering
            \import{tunnelbarrier/asym/irradiation/control}{Tbath_Afit.pgf}
            \caption{
                Measured bath temperature $T_\mathrm{bath}$ as a function of the effective fitted microwave amplitude $A_\mathrm{fit}$ at selected frequencies.
            }
            \label{fig:tb:asym:control-Tbath_Afit}
        \end{figure}

        \begin{figure}[t]
            \centering
            \import{tunnelbarrier/asym/irradiation/control}{Tfit_Tbath.pgf}
            \caption{
                Fitted transport temperature $T_\mathrm{fit}$ as a function of the measured bath temperature $T_\mathrm{bath}$ at selected frequencies. The clipped response separates the fitted-temperature floor, intermediate evolution, and upper saturation.
            }
            \label{fig:tb:asym:control-Tfit_Tbath}
        \end{figure}